
\begin{table}[htbp]
\centering
\small
\caption{Conceptual Synthesis of Identity and Its Related Ontological Dimensions}
\label{tab:identity-conceptual-synthesis}

\renewcommand{\arraystretch}{1.5}

\begin{tabular}{
    |p{0.19\textwidth}
    |p{0.32\textwidth}
    |p{0.35\textwidth}|
}
\hline

\textbf{Concept}
&
\textbf{Central Meaning}
&
\textbf{Relation to Identity}
\\
\hline

Identity
&
The ontological determinacy through which an existent is the existent that it is.
&
Foundational: inseparable from existence and not reducible to another concept.
\\
\hline

Determination
&
The determinacy of an existent as an identity-bearing existent.
&
Closely connected to identity, but not an additional operation that produces it.
\\
\hline

Distinction
&
The non-identity of an existent with what is other than it.
&
Articulates difference without exhausting or producing identity.
\\
\hline

Self-Sameness
&
The formal relation through which an existent is identical with itself.
&
Expresses a formal aspect of identity without providing its complete ontological explanation.
\\
\hline

Individuation
&
The particularity through which an existent is this existent rather than another.
&
Concerns particularity without being reducible to identity itself.
\\
\hline

Continuity
&
The problem of how an existent remains itself across successive conditions and transformations.
&
Concerns persistence and transformation without requiring absolute immutability.
\\
\hline

Ontological Relations
&
The connections and dependencies in which existents may participate.
&
May be relevant to the constitution or manifestation of identity, but cannot yet be treated as its exhaustive explanation.
\\
\hline

\end{tabular}

\end{table}
